\documentclass[11pt]{article}
\usepackage[utf8]{inputenc}
\usepackage[T1]{fontenc}
\usepackage{lmodern}
\usepackage[margin=1in]{geometry}
\usepackage{amsmath,amssymb}
\usepackage{natbib}
\usepackage{textcomp}
\usepackage[dvipsnames]{xcolor}
\usepackage{enumitem}
\usepackage{parskip}
\usepackage{titlesec}
\usepackage[colorlinks=true,linkcolor=blue,citecolor=blue,urlcolor=blue]{hyperref}

\titleformat{\section}{\large\bfseries}{\thesection}{1em}{}
\titleformat{\subsection}{\normalsize\bfseries}{\thesubsection}{1em}{}
\titleformat{\subsubsection}{\normalsize\bfseries\itshape}{}{0em}{}

\title{\textbf{Scale-Taxed, Not Scale-Invariant:\\ Measuring Whether Corpus-Level AI-Text Estimation Resists Budget-Bounded Evasion}}
\author{}
\date{}

\newcommand{\ph}{\textit{[placeholder --- fill in after experiments are run]}}

\begin{document}
\maketitle
\vspace{-2em}

\begin{quote}
\small\textit{Assembled from the six confirmed stage outputs of an \texttt{auto-research-discussion} run (seed topic B6.5), post research-taste audit (\textcolor{ForestGreen}{\textbf{PASS}}). Status: \textbf{specification / paper draft --- no experiments have been run.} All numbers below are pre-registered design targets, not results.}
\end{quote}

\begin{abstract}
Corpus-level estimators of AI-generated-text prevalence are increasingly treated as the safe alternative to single-document detection, precisely because instance-level verdicts are known to collapse under a bounded paraphrase. Whether that safety is real has not been directly measured: an averaging intuition --- that many independent per-document perturbations cancel out in an aggregate proportion --- is plausible but unverified, and the nearest existing measurement instrument validates its estimator only against a benign, non-adversarial mixing process. We formalize a cost-accuracy game between a budget-bounded paraphrase adversary and a corpus-level prevalence estimator, and ask whether the minimum budget required to bias the aggregate estimate by a fixed margin grows with corpus size, rather than only confirming robustness at one scale. We test whether a directionally coordinated adversary defeats the averaging argument that an independent one cannot; require the scaling conclusion to hold sign under two independent adversary cost models; validate the ground-truth construction with a shuffled-label negative control; and isolate whether any measured robustness is attributable to the estimator's specific aggregation design rather than to generic averaging over many documents. We report which of these hold for current corpus-level estimators and paraphrase adversaries, and what the result implies for treating population-scale AI-text auditing as strategy-proof.
\end{abstract}

\section{Introduction}

Synthetic, LLM-generated text is no longer a curiosity in web-scale corpora --- it is a routine contaminant that data-curation pipelines, content-authenticity audits, and platform-scale moderation systems all have to reckon with. A growing line of work responds by reframing AI-text detection as a measurement problem: rather than asking whether any single document was AI-written, it estimates what fraction of an entire corpus is. This reframe is attractive precisely because the alternative --- a verdict on each individual document --- is known to be fragile: a bounded, quality-preserving paraphrase is enough to defeat a state-of-the-art zero-shot detector, and that fragility is treated by this literature as settled ground it is explicitly built to route around.

Yet whether routing around single-document fragility actually buys anything has not been directly measured. A corpus-level estimator that computes only a mixing coefficient over aggregate token statistics does not obviously inherit the same vulnerability a per-document verdict has --- but it does not obviously avoid it either. An operator who assumes a pooled estimate is safe ``because any one verdict can be gamed but the average can't'' is trusting exactly the averaging intuition this paper sets out to test, not an established result. If that intuition is wrong, treating corpus-level auditing as a safe fallback is actively misleading. If it is right --- and specifically if the cost of defeating it grows with the size of the corpus being audited, rather than merely holding at whichever one scale happens to have been checked --- that is a positive, actionable finding: a genuinely different and better instrument, not single-document detection wearing an aggregation layer.

The nearest existing work approaches this gap from adjacent but distinct angles, and it is worth naming the closest of them plainly before surveying the rest: a recent reframing of AI-text detection as population-scale measurement --- the direct precedent for the corpus-level estimator this paper studies --- fits a mixture model over aggregate token statistics to estimate a corpus's AI-generated fraction, but validates that estimator only against a benign, non-adversarial mixing process; it never asks what happens when a paraphraser deliberately tries to move the fitted coefficient, which is precisely the gap this paper is built to close. A second line of work formalizes the vocabulary this paper's title borrows --- ``strategy-proof under a manipulation budget'' --- for incentive-aware bandit allocation, proving that confidence-based screening makes manipulation unprofitable; that is a result about arms in a bandit, however, with no existing bridge to a text corpus's aggregate statistic. Debiased proportion-estimation methods, in turn, correct per-record membership signals into an aggregate estimate under an assumed-stationary error model --- a template this paper borrows for a second estimator design --- but one built for ordinary measurement noise, not for an adversary who can shift the very error rates the correction relies on. And adaptive watermark-evasion work treats a budget-bounded attacker as the default threat model a defense must survive, not an edge case --- exactly the standard this corpus-level literature has not yet been held to --- though its per-text attack objective has never been pointed at a population-level target instead of a single verifier's verdict.

Turning this gap into a measurable claim surfaces its own difficulties. The entire case for corpus-level robustness rests on an unstated assumption that per-document paraphrase perturbations behave like independent noise that cancels under aggregation --- but a paraphraser is under no obligation to perturb independently, and whether a directionally coordinated rewrite allocation defeats that cancellation is a live, attacker-capability question distinct from whether the estimator's specific design matters at all. Reporting a cost in the first place is not innocent either: tokens rewritten, documents touched, and paraphrase compute imply different optimal attacks --- light uniform rewriting versus concentrated heavy rewriting of a subset --- so a scaling conclusion that only holds under one bookkeeping convention would be a property of the convention, not of the estimator. No real, uncurated corpus carries verifiable ground truth for what fraction of it is AI-generated, so any error metric first requires a construction that is hard enough to be informative while still admitting a checkable ground-truth proportion, without leaking a shortcut that lets an estimator ``succeed'' for the wrong reason. A debiased estimator's correction is only as good as the stationarity of the error rates it was calibrated against, and an adversary that shifts those rates by paraphrasing breaks that calibration in a way a likelihood-based estimator's failure mode does not resemble. And even a demonstrated robustness advantage says nothing about \emph{why} it exists unless it is checked against the null hypothesis that any sufficiently large average would show the same advantage, dedicated estimator or not.

Our approach turns each of these obstacles into a specific design choice rather than a hedge. A coordinated, adaptively-trained adversary is pitted against the same non-adaptive paraphraser at matched budget, isolating whether the estimator's apparent robustness depends on the weaker assumption of an uncoordinated attacker. The same scaling measurement is repeated under two independently defined, non-nested cost models, and only a sign-consistent result across both is reported as a property of the estimator rather than of one bookkeeping convention. Ground truth is constructed from matched human/LLM text pairs and checked against a shuffled-label negative control before any robustness claim is interpreted at all. A debias-and-aggregate estimator is run under the same adversary that targets its stationarity assumption directly, alongside the primary likelihood-based estimator, to see which design's specific failure mode actually manifests. And a naive average of per-document detector scores --- no calibration, no distributional modeling --- is run under identical attack conditions as a null model, so that any measured robustness can be attributed to estimator design rather than to the generic fact that averages of large populations have lower variance than single draws.

We evaluate this design on semi-synthetic corpora built from real human and LLM-generated text spanning multiple generator families, at corpus sizes sweeping several orders of magnitude, comparing a population-scale single-document-mean control and a budget-matched coordinated adversary against the dedicated corpus-level estimators.

In summary, we contribute:
\begin{itemize}[leftmargin=1.4em]
\item A cost--accuracy formalization of ``evasion economics'' for corpus-level AI-text prevalence estimation, giving the intuitive claim that aggregation resists evasion a precise, falsifiable form: a scaling exponent relating the minimum bias-inducing adversary budget to corpus size, rather than a single-scale accuracy number.
\item A coordinated-adversary isolation that tests whether corpus-level robustness depends on the weaker assumption that a paraphraser perturbs documents independently, together with a direct check of whether the scaling property survives when that assumption is dropped.
\item A cost-model-invariance requirement and a shuffled-label ground-truth validity check, together ruling out two distinct ways the central scaling claim could be an artifact of the measurement setup rather than a property of the phenomenon.
\item A naive-averaging null model that isolates whether any observed robustness is attributable to the corpus estimator's specific aggregation design, as opposed to the generic variance reduction of averaging over a large population.
\item An empirical account of whether, and under what threat-model conditions, population-scale AI-text auditing is more strategy-proof than the single-document detection it is meant to replace.
\end{itemize}

\section{Problem Definition}

\noindent\textbf{Formal setup.} Consider a corpus $D_N = \{d_1, \dots, d_N\}$, where each document $d_i$ carries a latent label $y_i \in \{\text{human}, \text{AI}\}$, unknown to any estimator but controllable by construction in a semi-synthetic corpus. The true population-level AI-generation proportion is $\alpha = \frac{1}{N}\sum_i \mathbb{1}[y_i = \text{AI}]$. A corpus-level estimator $E$ returns $\hat\alpha = E(D_N)$, and a single-document detector $g$ returns a per-document score whose corpus-wide discriminative power is summarized by $\mathrm{AUROC}(g; D_N)$. An adversary $\mathcal{A}$ transforms the corpus $D_N \to D_N'$ by paraphrasing a subset of documents, spending a total budget $B(\mathcal{A}) = \sum_i c(d_i, d_i')$ under a declared cost model $c$ (per token rewritten, or per document touched scaled by rewrite strength). Write $e_0(N) = \mathbb{E}\,|E(D_N) - \alpha|$ for the estimator's clean (no-adversary) error and $e_B(N) = \mathbb{E}\,|E(D_N') - \alpha|$ for its error after an attack of budget $B$. The bias margin induced by an attack is $\varepsilon = |E(D_N') - E(D_N)|$, and the minimum budget needed to induce a bias of at least $\varepsilon$ is $B_{\min}(\varepsilon, N) = \min\{B(\mathcal{A}) : |E(D_N') - E(D_N)| \ge \varepsilon\}$. The object of study is the joint behavior of $e_B(N)$ and of the scaling exponent $\beta$ --- the ordinary-least-squares slope of $\log B_{\min}(\varepsilon, N)$ on $\log N$ --- relative to $g$'s collapse under the identical budget $B$.

\noindent\textbf{Why this is hard.} The case for $e_B(N)$ staying bounded rests on an implicit assumption that per-document perturbations behave like independent noise that averages out in $\hat\alpha$; nothing about the estimator's definition rules out a paraphraser that instead perturbs documents in a coordinated, directional way specifically chosen to move $\hat\alpha$, and whether that assumption is load-bearing is exactly what has to be checked rather than presumed. The budget $B(\mathcal{A})$ itself is only meaningful relative to a declared cost model $c$, and two natural, non-nested choices --- cost proportional to tokens rewritten versus cost proportional to documents touched --- reward different optimal attacks (uniform light rewriting versus concentrated heavy rewriting of a subset), so a conclusion about the sign of $\beta$ that holds under only one $c$ is a property of that convention rather than of $E$. Because $\alpha$ is unobservable in any real, uncurated corpus, computing $e_0(N)$ or $e_B(N)$ at all requires a semi-synthetic construction with an experimenter-set ground truth, and that construction must be checked against the possibility that some artifact of its assembly --- rather than the estimator itself --- is what makes $\alpha$ recoverable. Finally, a demonstrated gap between $e_B(N)$ and a naive per-document-average baseline's error under the same attack says nothing about \emph{why} the gap exists unless the two are compared directly at matched $(N, B)$ --- a large-$N$ average could show reduced variance under attack for reasons that have nothing to do with $E$'s particular aggregation design.

\section{Method Overview}

At a conceptual level, the study is a single cost--accuracy measurement pipeline built around one estimator-vs-adversary pairing, with four isolation controls layered on top of it to rule out the specific ways the resulting scaling claim could be an artifact rather than a property of the phenomenon. A corpus-level estimator and a budget-bounded paraphrase adversary enter the pipeline together; the adversary's minimum budget to bias the estimate by a fixed margin is measured across a sweep of corpus sizes, yielding the scaling exponent $\beta$ that is this paper's central quantity. Each control then attacks one specific way that measurement could be misleading: whether the adversary's \emph{coordination} (not just its budget) matters, whether the \emph{cost model} used to report the budget changes the conclusion's sign, whether the \emph{ground-truth construction} itself is trustworthy, and whether the observed robustness is attributable to the estimator's \emph{specific design} rather than to averaging in general.

The \textbf{coordinated adversary arm} takes the same budget-bounded paraphrase adversary and adds a directionally-optimized variant, trained against a surrogate of the estimator, to test whether the independence assumption behind the base measurement is load-bearing. The \textbf{cost-model invariance check} re-runs the entire minimum-budget search under a second, independently defined cost model and requires the sign of the fitted scaling exponent to agree with the first. The \textbf{semi-synthetic ground-truth design} builds the evaluation corpora themselves, from matched human/LLM text pairs, and validates them with a shuffled-label negative control before any downstream robustness number is interpreted --- and, on that validated instrument, delivers the paper's headline attacked-versus-clean comparison. The \textbf{debiasing-under-shift diagnostic} runs a second, debias-and-aggregate estimator design under the same adversary and checks whether its stationarity assumption fails in a way the primary estimator's does not. The \textbf{naive-average confound control} runs an estimator with no distributional design at all --- a plain average of per-document detector scores --- under identical attack conditions, so that any robustness gap can be attributed to the dedicated estimators' design rather than to generic large-sample averaging.

\section{Methodology}

\subsection{Coordinated Adversary Arm}
\textbf{Challenge.} The case for the estimator's robustness rests on an unstated assumption that per-document paraphrase perturbations behave like independent noise that cancels under aggregation; nothing rules out a paraphraser that instead coordinates its rewrites directionally to move the aggregate estimate.

\textbf{Insight.} Independent per-document noise cancels under aggregation by construction, but a directionally coordinated adversary specifically breaks that cancellation --- so pitting a coordinated adversary against a non-adaptive one at matched budget isolates whether the estimator's apparent robustness depends on the (weaker, and not obviously realistic) assumption of an uncoordinated attacker.

\textbf{Method.} Alongside a non-adaptive, fixed-budget paraphraser that rewrites documents independently, a second, adaptively-trained adversary is optimized against a surrogate of the estimator to bias the aggregate estimate as efficiently as possible within the same declared budget. The two adversaries' budget-to-bias efficiency is compared at matched corpus size, and --- if the coordinated adversary requires strictly less budget --- the scaling measurement is repeated under it specifically, to check whether the scale-taxed cost structure survives even the stronger of the two attackers considered.

\subsection{Cost-Model Invariance Check}
\textbf{Challenge.} An adversary's ``budget'' is only meaningful relative to a declared cost model, and two natural conventions --- cost proportional to tokens rewritten versus cost proportional to documents touched --- reward different optimal attacks, so a scaling conclusion computed under only one convention risks being an artifact of that convention.

\textbf{Insight.} If the qualitative conclusion (the sign of the fitted scaling exponent) is required to hold under two independently defined, non-nested cost models rather than reported under whichever one is more favorable, the conclusion is checked to be a property of the estimator rather than of the bookkeeping.

\textbf{Method.} The minimum-budget-to-bias search is run twice per corpus size --- once under a per-token-rewritten cost model, once under a per-document-touched cost model scaled by rewrite strength --- and the fitted scaling exponent's sign is compared across the two. Only a sign-consistent result across both is carried forward as a claim about the estimator; a sign disagreement is reported as cost-model-dependent rather than resolved by preferring one convention.

\subsection{Semi-Synthetic Ground-Truth Design}
\textbf{Challenge.} No real, uncurated corpus has verifiable ground truth for what fraction of it is AI-generated, so computing any error metric at all requires a construction that is informative without leaking a shortcut that lets an estimator ``succeed'' for the wrong reason.

\textbf{Insight.} A ground-truth check that only confirms an estimator recovers the intended proportion is not enough on its own; a shuffled-label negative control is what rules out the possibility that the construction procedure itself, rather than the estimator, is responsible for an accurate recovery --- and only once that instrument is validated does the paper's headline attacked-versus-clean comparison become interpretable.

\textbf{Method.} Semi-synthetic corpora are built from matched human/LLM text pairs with an experimenter-set ground-truth proportion $\alpha$. Two checks are run before any downstream claim: the clean (no-adversary) estimator must recover $\alpha$ within a pre-registered tolerance, and a shuffled-label control --- in which the human/AI labels are permuted independently of document content --- must fail to recover $\alpha$ by a significant margin. On that validated construction, the headline comparison is then run: under a fixed adversary budget that drives the paired single-document detector's AUROC below a fragility threshold, the corpus-level estimator's attacked error is compared against its clean-error baseline.

\subsection{Debiasing-Under-Shift Diagnostic}
\textbf{Challenge.} A debias-and-aggregate estimator corrects per-document detector scores using reference error rates assumed to be stationary; an adversary that paraphrases specifically to shift those per-document error rates violates the very assumption the correction depends on, a distinct failure mode from a likelihood-based estimator's.

\textbf{Insight.} A debiasing correction computed once from clean reference data is only valid for as long as the adversary leaves the error rates it corrects for unchanged --- checking this directly, rather than assuming it, separates a stationarity-violation failure from the coordination question addressed by the coordinated adversary arm above.

\textbf{Method.} A debias-and-aggregate estimator is run alongside the primary likelihood-based estimator under the identical fixed-budget adversary, and each estimator's attacked error is compared against its own no-adversary baseline at matched budget and corpus size, isolating whether the debiasing correction's stationarity assumption is the more fragile of the two estimator designs.

\subsection{Naive-Average Confound Control}
\textbf{Challenge.} A demonstrated gap between a dedicated corpus estimator's attacked error and a naive baseline's says nothing about \emph{why} the gap exists unless the two are compared directly --- a large sample's average could show reduced variance under attack for reasons unrelated to any specific aggregation design.

\textbf{Insight.} If a naive average of per-document detector scores, with no calibration or distributional modeling, is just as robust as the dedicated estimators under identical attack conditions, the observed ``corpus-level robustness'' is a generic consequence of averaging over many documents, not evidence that a specifically designed estimator does anything --- making this comparison what turns the estimator-design contribution into a testable claim rather than an assumed one.

\textbf{Method.} A naive estimator that simply averages the single-document detector's per-document scores into a corpus-level proportion, with no calibration or distributional modeling, is run under identical attack conditions to the two dedicated estimators at matched corpus size and budget, and all three estimators' attacked errors are compared directly.

\section{Experiments}

\subsection{Experimental Setup}

\textbf{Datasets.} RAID \citep{RAID2024} --- a large multi-domain, multi-generator machine-generated-text benchmark including adversarial/paraphrase variants; supplies at least three generator families (GPT-, Llama-, Claude-class) so the scaling result is not an artifact of one generator. HC3 \citep{HC32023} --- paired human/LLM answers to the same prompts, giving the tightest available control for constructing controllable-$\alpha$ semi-synthetic corpora with minimal topical confound between the human and AI portions; this directly stresses the shuffled-label control, since topically-matched pairs make a topic-correlated shortcut a real risk to rule out rather than one simulated away by construction. Semi-synthetic mixed corpora with an experimenter-controlled ground-truth AI proportion, built from the two sources above, following the semi-synthetic validation protocol of the nearest prior corpus-level estimator \citep{MonitorAIScale}.

\textbf{Baselines} (sanity-check comparators, not competitors to beat --- this is a measurement study, and no research question here is framed as ``our procedure beats baseline X''). A population-scale single-document-mean control --- a plain average of per-document detector scores, with no calibration or distributional modeling --- is the ``why do this at all'' check: if this control is already strategy-proof, the entire dedicated-estimator apparatus is unjustified. A budget- and capability-matched coordinated adversary --- the same instrument as the coordinated adversary arm above --- is the ``strongest-alternative-explanation'' check, ruling out that any observed robustness is merely an artifact of a weak, non-adaptive attacker.

\textbf{Metrics.} The corpus estimator's attacked and clean error, $e_B(N)$ and $e_0(N)$; the paired single-document detector's AUROC before and after the identical attack; the minimum adversary budget $B_{\min}(\varepsilon, N)$ to bias the estimate by a fixed margin; the fitted scaling exponent $\hat\beta$ of $B_{\min}$ against corpus size $N$, by log-log regression with bootstrap confidence intervals.

\textbf{Compute plan (budget-search amortization).} The $B_{\min}$ search is the design's dominant cost, and run naively it is infeasible: the RQ3 grid alone is 4 corpus sizes $\times$ 2 cost models $\times$ 3--5 seeds $\approx$ 40 cells, each requiring a search over budgets ($\approx$8--10 evaluations under bisection), and each evaluation re-paraphrasing up to $O(N)$ documents --- at $N = 10^5$, on the order of $10^6$--$10^7$ fresh LLM paraphrase calls per cell. The design therefore amortizes generation: a \emph{paraphrase ladder} --- a small fixed set of rewrite variants at graded strengths --- is pre-generated exactly once per document per generator family, and every subsequent budget-search step, cost model, adversary arm (non-adaptive and coordinated), and seed \emph{selects among cached variants} rather than invoking the paraphraser, making the marginal cost of a $B_{\min}$ evaluation selection and re-scoring, not generation. Bisection over the budget axis is licensed by an explicitly stated monotonicity assumption --- that the maximum bias achievable by the optimal allocation is non-decreasing in the budget, which holds by construction since any lower-budget allocation remains feasible at a higher budget. Cost declared per cost model is computed over the \emph{selected} ladder rungs, so the cached-generation shortcut changes what is computed nowhere --- only when the paraphrases are produced.

\textbf{Research questions:}
\begin{enumerate}[leftmargin=1.6em]
\item RQ1 --- Does a coordinated, adaptively-optimized adversary require no less budget than a non-adaptive one to bias the estimate by the same margin?
\item RQ2 --- If a coordinated adversary does reduce the required budget, does the scaling exponent's confidence interval still exclude zero?
\item RQ3 --- Does the scaling exponent's confidence interval exclude zero under both declared cost models?
\item RQ4 --- Does the clean estimator recover ground truth within tolerance while a shuffled-label control fails to?
\item RQ5 --- Does the corpus estimator's error stay bounded under a budget that already collapses the paired single-document detector?
\item RQ6 --- Does a debias-and-aggregate estimator's error grow more than a likelihood-based estimator's under the identical adversary?
\item RQ7 --- Does a naive per-document-average estimator's error exceed each dedicated estimator's by a substantial margin under the same attack?
\end{enumerate}

\subsubsection{RQ1: Does a coordinated, adaptively-optimized adversary require no less budget than a non-adaptive one to bias the estimate by the same margin?}
\textbf{Setup.} At a fixed corpus size, train the coordinated adversary against a surrogate of the corpus estimator via preference optimization on best-of-N rewrite candidates (no gradient access to the true estimator assumed), then measure the minimum budget each of the coordinated and non-adaptive adversaries requires to bias the estimate by margin $\varepsilon = 0.10$, at matched corpus size and cost model. Sample size: 3--5 seeds per adversary type, repeated across the same generator families used elsewhere in the grid.

\textbf{Expected outcome.} The coordinated adversary requires no less budget than the non-adaptive adversary to reach the same bias margin --- the aggregation-based robustness advantage does not depend on the adversary being uncoordinated.

\textbf{Falsification condition.} Refuted if the coordinated adversary requires strictly less budget than the non-adaptive adversary at matched margin and corpus size --- that would mean the base robustness measurement depends on an unrealistically weak (independent-perturbation) threat model, and the finding would need to be qualified to that narrower model rather than stated unconditionally.

\textbf{Hypothesis (plain language).} Coordinating rewrites, rather than applying them independently, does not by itself make evasion of the aggregate estimate cheaper.

\textbf{Metrics.} Minimum budget to reach $\varepsilon = 0.10$ for each adversary type; ratio of coordinated-to-non-adaptive minimum budget with bootstrap confidence interval.

\textbf{Result.} \ph

\subsubsection{RQ2: If a coordinated adversary does reduce the required budget, does the scaling exponent's confidence interval still exclude zero?}
\textbf{Setup.} Conditional on RQ1 finding a budget reduction under coordination, repeat the full corpus-size sweep and log-log regression of minimum budget against $N$ using the coordinated adversary in place of the non-adaptive one, at the primary cost model.

\textbf{Expected outcome.} The fitted scaling exponent $\hat\beta$'s 95\% bootstrap confidence interval lower bound remains greater than 0 under the coordinated adversary --- coordination may narrow the scaling advantage but does not eliminate it.

\textbf{Falsification condition.} Refuted if the confidence interval on $\hat\beta$ under the coordinated adversary includes or falls below 0 --- that would mean the scale-taxed cost structure is specific to weak, non-adaptive attackers and does not hold against a realistic, budget-matched competent one.

\textbf{Hypothesis (plain language).} Even a stronger, coordinated adversary still faces a rising cost to bias the aggregate estimate as corpus size grows.

\textbf{Metrics.} $\hat\beta$ under the coordinated adversary, with bootstrap 95\% confidence interval, compared against the base non-adaptive-adversary $\hat\beta$ fitted at the primary cost model inside RQ3's corpus-size sweep (RQ3's question is cost-model invariance; its sweep is where the base non-adaptive scaling exponent is computed).

\textbf{Result.} \ph

\subsubsection{RQ3: Does the scaling exponent's confidence interval exclude zero under both declared cost models?}
\textbf{Setup.} Sweep corpus size $N \in \{10^2, 10^3, 10^4, 10^5\}$ under the non-adaptive adversary, fitting the minimum budget to bias $\varepsilon = 0.10$ separately under a per-token-rewritten cost model and a per-document-touched cost model, 3--5 seeds per (N, cost model) cell.

\textbf{Expected outcome.} The fitted scaling exponent $\hat\beta$ has a 95\% bootstrap confidence interval lower bound greater than 0, with a point estimate at or above 0.5, under both cost models.

\textbf{Falsification condition.} Refuted (no scaling advantage) if either cost model's confidence interval on $\hat\beta$ includes or falls below 0; refuted as cost-model-dependent (a weaker, qualified finding) if the two cost models disagree in sign even when each individually excludes 0.

\textbf{Hypothesis (plain language).} The minimum budget to bias the corpus-level estimate grows with corpus size regardless of which reasonable way of counting adversary cost is used.

\textbf{Metrics.} $\hat\beta$ per cost model with bootstrap 95\% confidence interval; sign agreement across cost models.

\textbf{Result.} \ph

\subsubsection{RQ4: Does the clean estimator recover ground truth within tolerance while a shuffled-label control fails to?}
\textbf{Setup.} On each semi-synthetic corpus (no adversary present), run the clean corpus estimator and compare its recovered $\hat\alpha$ against the experimenter-set $\alpha$; separately, permute the human/AI labels independently of document content and rerun the identical estimator on this shuffled-label corpus.

\textbf{Expected outcome.} The clean estimator recovers $\alpha$ within a pre-registered tolerance across all tested corpus sizes, while the shuffled-label control's estimate deviates from $\alpha$ by significantly more than that tolerance.

\textbf{Falsification condition.} Refuted if the shuffled-label control also recovers $\alpha$ within the same tolerance --- that would indicate the semi-synthetic construction leaks a shortcut (for instance, a topic or style correlate) that lets any estimator appear accurate regardless of whether it is measuring authorship at all, invalidating every downstream robustness claim built on this construction.

\textbf{Hypothesis (plain language).} The ground-truth construction is measuring what it claims to measure, not an incidental artifact of how the corpus was assembled.

\textbf{Metrics.} $|\hat\alpha - \alpha|$ for the clean and shuffled-label conditions, compared against the pre-registered tolerance.

\textbf{Result.} \ph

\subsubsection{RQ5: Does the corpus estimator's error stay bounded under a budget that already collapses the paired single-document detector?}
\textbf{Setup.} At $N = 10^5$, apply the non-adaptive paraphrase adversary at a budget calibrated to drive the paired single-document detector's AUROC to at most 0.60; measure the corpus estimator's error under that same budget and compare it against its clean-error baseline. Sample size: 3--5 seeds $\times$ at least 3 generator families, Holm-corrected across generators.

\textbf{Expected outcome.} The corpus-level estimator's attacked error $e_B(10^5)$ remains at most 0.05 and at most twice its clean error $e_0(10^5)$, while the paired single-document detector's AUROC falls to at most 0.60 under the identical budget.

\textbf{Falsification condition.} Refuted if $e_B(10^5)$ exceeds twice $e_0(10^5)$ (or exceeds 0.05 outright) at a budget that collapses the paired detector's AUROC to 0.60 or below --- that would mean corpus-level estimation inherits single-document fragility rather than resisting the same attack budget that defeats it.

\textbf{Hypothesis (plain language).} Corpus-level estimation resists the same adversary budget that reliably breaks single-document detection.

\textbf{Metrics.} $e_B(N)$, $e_0(N)$, and $\mathrm{AUROC}(g; D_N')$ at $N = 10^5$, with bootstrap 95\% confidence intervals.

\textbf{Result.} \ph

\subsubsection{RQ6: Does a debias-and-aggregate estimator's error grow more than a likelihood-based estimator's under the identical adversary?}
\textbf{Setup.} At matched corpus size and adversary budget, run both the primary likelihood-based (mixture-MLE) estimator and a debias-and-aggregate estimator under the identical non-adaptive adversary, comparing each estimator's attacked error against its own clean-error baseline.

\textbf{Expected outcome.} The debias-and-aggregate estimator's error grows relative to its no-adversary baseline by more than the likelihood-based estimator's does, at matched budget and corpus size --- the stationarity assumption behind debiasing is the more fragile of the two estimator designs under this attack.

\textbf{Falsification condition.} Refuted if the two estimators' relative error growth under attack is statistically indistinguishable, or if the debias-and-aggregate estimator is in fact the more robust of the two --- either result would mean the stationarity-violation failure mode this diagnostic targets is not actually present, or not the dominant one, at the tested budgets.

\textbf{Hypothesis (plain language).} An estimator that debiases using an assumed-stationary error correction is more exposed to an adversary that can shift that error rate than a likelihood-based estimator is.

\textbf{Metrics.} Relative error growth $e_B(N)/e_0(N)$ for each estimator design, with bootstrap 95\% confidence interval on the difference.

\textbf{Result.} \ph

\subsubsection{RQ7: Does a naive per-document-average estimator's error exceed each dedicated estimator's by a substantial margin under the same attack?}
\textbf{Setup.} At matched corpus size and adversary budget, run the naive per-document-average estimator alongside both dedicated corpus estimators (mixture-MLE, debias-aggregate) under identical attack conditions.

\textbf{Expected outcome.} The naive estimator's attacked error exceeds each dedicated estimator's by at least $2\times$, with the ratio's 95\% bootstrap confidence interval lower bound at or above 2 --- indicating the robustness benefit is attributable to estimator design, not to generic averaging over a large population.

\textbf{Falsification condition.} Refuted if the naive estimator's error is statistically indistinguishable from (within $1.25\times$ of) either dedicated estimator's --- that would mean any observed corpus-level robustness is a generic consequence of averaging over many documents rather than evidence that a specifically designed estimator does anything, reframing the contribution from ``estimator design matters'' to ``aggregation alone, of any kind, is what matters.''

\textbf{Hypothesis (plain language).} A dedicated estimator's specific aggregation design, not just the fact of averaging over many documents, is doing real work.

\textbf{Metrics.} $e_B(N)$ for the naive, mixture-MLE, and debias-aggregate estimators at matched $(N, B)$; ratio of naive to each dedicated estimator's error, with bootstrap confidence interval.

\textbf{Result.} \ph

\vspace{0.5em}
\noindent\textbf{Further evaluation.} With more time, a natural extension is to test a fully adaptive, estimator-omniscient (white-box) adversary that directly targets the exact aggregation statistic, rather than the budget-bounded, non-omniscient adversary this design studies throughout --- explicitly out of scope here because it changes the threat model from ``bounded, surrogate-optimized attacker'' to ``attacker with full knowledge of the defense,'' which would need its own identifiability design rather than an extension of the coordinated-adversary arm above.

\section{Related Work}

The closest prior work to this paper is the corpus-level measurement reframing that first estimated AI-text prevalence by fitting a mixture model over aggregate token statistics \citep{MonitorAIScale}, and it is worth stating the delta plainly before surveying the wider landscape: that estimator is the direct precedent this paper's own primary corpus-level instrument is built on, but it validates its estimator only against a benign, non-adversarial mixing process and never asks what happens when a paraphraser deliberately tries to move the fitted coefficient --- precisely the identifiability gap this paper is built to close. A second, more distant anchor is the strategy-proofness literature for incentive-aware bandit allocation \citep{UCBS}, together with a peer-prediction mechanism-design literature that adaptively learns incentive-compatible rewards for strategic agents under no prior on their beliefs \citep{MultiAgentMechDesign2025}; both supply the vocabulary this paper's own title borrows --- ``robust to manipulation under a budget'' --- via confidence-based screening or online-learned payments that provably make manipulation unprofitable, but for arms in a bandit or reports from a rational agent, with no existing bridge to a text corpus's aggregate statistic.

\textbf{Corpus-level and population-scale AI-text measurement.} The mixture-model reframing above estimates an AI-generated fraction by maximum likelihood over aggregate token-occurrence statistics, validated on semi-synthetic mixtures with known ground-truth proportions \citep{MonitorAIScale}; a companion line of work on debiased proportion estimation corrects noisy per-record membership-inference signals into an unbiased aggregate estimate with asymptotic confidence intervals \citep{DUCI}, a general debias-and-aggregate template this paper adapts as a second corpus-level estimator design. Both establish that a population-level proportion can be estimated accurately under benign conditions, but neither models an adversary who paraphrases, or otherwise manipulates per-record signals, specifically to bias the aggregate --- the gap this paper's adversary-vs-estimator game is built to fill. A third, purely theoretical line of work derives matching minimax upper and lower bounds on the sample complexity of estimating a corpus-level watermark proportion from Gumbel-Max-watermarked text, under both full and pivotal-statistic observation regimes \citep{GumbelWatermarkSampleComplexity2026}. That result establishes exactly the estimation-theoretic difficulty this paper's own mixture-MLE estimator navigates empirically --- how many samples are needed to pin down a population-level proportion at all --- but it is a sample-complexity result under a purely benign, non-adversarial sampling process; it says nothing about what happens once an adversary is allowed to shift the very statistics the estimator is built to recover, which is precisely this paper's question.

\textbf{Single-document detection and its known fragility.} Zero-shot detectors such as the log-probability-curvature approach that compares a passage's likelihood to that of small semantics-preserving perturbations produce a per-document verdict whose accuracy is known to collapse under paraphrasing \citep{DetectGPT}; detector-agnostic calibration layers that refine per-token scores via pairwise mean-field inference before aggregating to a passage-level decision improve robustness to benign sampling noise but still target a single-passage decision, not a corpus proportion, and assume noise is incidental rather than adversarially coordinated \citep{MarkovCalib}. Both establish the per-document fragility this paper takes as a settled premise and retains as a paired control, but neither offers, nor claims to offer, an aggregate-estimate quantity that a corpus-level robustness claim could be checked against. The paraphrase-fragility premise itself now rests on a mature attack literature: recursive paraphrasing defeats detectors and watermarks alike, with retrieval over prior generations the strongest known defense \citep{KrishnaParaphrase2023}; theoretical arguments bound any detector's advantage by the distributional distance a paraphraser can shrink \citep{Sadasivan2023}; and the current attack frontier is detector-guided --- an off-the-shelf LLM paraphrases under the guidance of a reference detector \citep{AdvParaphrase2025} --- with dedicated benchmarks standardizing paraphrase-attack evaluation \citep{PADBen2025} and black-box style-humanization attacks extending it \citep{MASH2026}. Every one of these attacks, and every benchmark scoring them, defines success against a per-document verdict; none measures, or is instrumented to measure, displacement of a corpus-level prevalence estimate --- which is exactly the attack surface this paper measures.

\textbf{Quantification learning.} A reader from the quantification-learning community would recognize this paper's three estimator designs as instances of that field's standard families: the naive per-document-score average is (probabilistic) classify-and-count, the debias-and-aggregate design is an adjusted-count correction, and the mixture-MLE estimator is a distribution-matching quantifier --- a taxonomy developed over two decades of work on class-prevalence estimation under prior-probability shift, with variance-minimizing estimator design \citep{QuantifierVariance2021} and systematic empirical comparison of the families under dataset shift \citep{QuantShift2023}. This paper claims no novelty for the estimator families themselves, and adopts that field's central lesson --- that aggregation design, not classifier accuracy, drives prevalence-estimation error. What the quantification literature does not contain, to our knowledge, is a strategic adversary: its distribution shift is nature-given, arising between training and deployment, whereas the shift studied here is chosen by a budget-bounded optimizer specifically to bias the quantifier. That adversarial-quantification question appears unposed in the field, and its answer --- how the manipulation budget must scale with population size --- is precisely the exponent this paper measures.

\textbf{Budget-bounded adversarial evasion.} Adaptive watermark-evasion work formalizes evasion as an explicit optimization problem --- expected verification failure traded against a quality-preservation term --- and trains a low-cost, budget-constrained paraphrase attacker via preference optimization against a surrogate verifier, arguing that adaptive attackers, not static ones, should be the default threat model for any detection-adjacent claim \citep{AdaptiveWMAttack}. This paper adapts that adversary-construction protocol directly for its own coordinated-adversary arm, but the original work's attack objective and budget are defined per text against a per-text verifier; it never asks whether the same per-document budget suffices to move a population-level statistic instead of a single verdict, which is exactly the question this paper's cost-model-invariant scaling measurement is built to answer. A mirror-image line of work rewrites adversarially perturbed inputs to restore a downstream classifier's correct predictions, without retraining the classifier itself \citep{ATINTER2023} --- the same rewrite-as-intervention mechanism this paper's adversary wields offensively against an estimator, deployed there defensively for ordinary text classifiers, with no treatment of a corpus-level aggregate statistic.

\textbf{Instance-level detection benchmarks that cite population-scale claims.} Recent peer-review-integrity work builds a manuscript-conditioned instance-level detector benchmarked in the strict low-false-positive regime, explicitly motivated by prior population-level estimates of LLM involvement in reviews \citep{AnchorPeerReview}. This work demonstrates that instance-level detection and population-level prevalence claims coexist in the same applied literature, but it neither analyzes the robustness of the population-level estimates its own motivation cites, nor connects its low-false-positive instance layer to an aggregate estimator's error under attack --- leaving exactly the connection this paper's cost--accuracy game is designed to make explicit and falsifiable.

\section{Threats to Validity}

\textbf{External.} All findings are demonstrated on synthetic/semi-synthetic corpora built from a limited set of LLM generator families and domains (RAID/HC3-derived text) and may not generalize to the full diversity of real-world contaminated corpora or to unseen generator models. \emph{Mitigation:} at least three generator families (GPT-, Llama-, Claude-class) and multiple RAID domains are included by design, with per-generator and per-domain breakdowns reported rather than a single pooled number.

\textbf{Construct.} The adversary ``budget'' (tokens rewritten, or documents touched scaled by rewrite strength) is a simplification of real adversarial cost, which also involves detection-of-evasion risk, latency, and coordination overhead not captured by a scalar quantity. More fundamentally, both cost models are \emph{declared bookkeeping conventions with no external anchor}: neither is calibrated against a market price of paraphrase compute, an API rate card, or any observed attacker's cost function, because no such ground truth is accessible --- the entire cost economy in this study is simulated. \emph{Mitigation (and its honest limit):} the scaling result is reported under two independent, non-nested cost models (the Cost-Model Invariance Check), and the strategy-proofness conclusion's sign is required to agree across both before being reported as a property of the estimator --- but sign-invariance across two conventions establishes invariance \emph{within the declared cost family only}. The measured exponent $\beta$ is a property of the estimator under that declared family; no claim in this paper should be read as a statement about real-world attacker economics in currency units, and the phrase ``economics of evasion'' throughout refers to this declared cost geometry, not to market costs.

\textbf{Internal.} The core comparison (corpus estimator vs.\ single-document detector) may share the same underlying detection signal, so any robustness advantage could be an artifact of the specific aggregation method rather than a general property of corpus-level estimation. \emph{Mitigation:} the Naive-Average Confound Control (RQ7) directly tests this by comparing the naive per-document average against two independent estimator families (mixture-MLE and debias-and-aggregate) under identical attack conditions.

\textbf{Conclusion validity.} The scaling exponent $\beta$ is estimated from a small range of corpus sizes $N$ with limited seeds, so an apparent budget-grows-with-$N$ trend could be statistical noise or a small-$N$ artifact. \emph{Mitigation:} $\beta$ is fit by log-log regression with bootstrap 95\% confidence intervals against a pre-registered one-sided test $H_0: \beta \le 0$, and the confidence interval must exclude 0 before any scaling effect is claimed (RQ3's falsification condition).

\textbf{Scope.} The threat model throughout is a budget-bounded, non-omniscient adversary; a fully adaptive, white-box adversary that targets the exact aggregation statistic directly is explicitly out of scope (see ``Further evaluation'' above) and could behave differently --- this is a declared boundary of the claim, not an oversight, and no result in this design should be read as a claim about that stronger threat model.

\section{Conclusion}

This work reframes an assumption implicit in a growing measurement practice --- that pooling AI-text judgments over a whole corpus is inherently safer than judging any one document --- as an empirical question with a precise, falsifiable form: not ``is the corpus estimate accurate today,'' but ``does the cost of biasing it grow with the size of the corpus being audited.'' The central methodological wager is that this question cannot be answered by a single accuracy number at one corpus size; it requires a scaling exponent, checked for sign agreement across two independent cost models, and checked against a coordinated (not merely budget-bounded) adversary, a naive-averaging null model, and a shuffled-label ground-truth control --- weaker designs risk reporting an artifact of the cost model, the attacker, or the construction as if it were a property of the estimator. Whichever direction the measured exponent points, the finding is decision-relevant: a positive, sign-stable scaling exponent that survives a coordinated attacker would mean population-scale AI-text auditing is a genuinely different and more trustworthy instrument than single-document detection, not a repackaging of it; a null or negative exponent, or one that depends on which adversary or cost model is used, would mean treating corpus-level estimates as inherently safe is premature and possibly dangerous. Three limitations bound this claim from the outset and are natural next steps rather than afterthoughts: the adversary throughout is explicitly budget-bounded and non-omniscient, leaving a fully adaptive, estimator-aware attacker for future work rather than folding it into the present threat model; the entire study is conducted on semi-synthetic corpora built from a limited set of generator families and domains, so generalization to the full diversity of real-world contaminated corpora and unseen generator models remains an open question this design reports rather than resolves; and the adversary's cost economy is simulated end to end --- both cost models are declared conventions with no external anchor in observed attacker behavior or market prices, so the ``economics'' this paper measures is the cost geometry the declared family induces, a scoping that an eventual deployment-grounded study (with real attacker cost data) would need to close before any operational budget threshold could be read off the fitted exponent.

\vspace{1em}
\hrule
\vspace{1em}

\section*{Pre-handoff checklist}
\begin{itemize}[leftmargin=1.4em]
\item \textbf{Organization/navigation} --- narrative flows problem $\to$ method $\to$ experiments $\to$ related work without forward references to undefined terms.
\item \textbf{Terminology discipline} --- ``cost model,'' ``scaling exponent,'' ``bias margin,'' ``coordinated adversary'' each defined at first use in Problem Definition, used consistently thereafter; no internal stage IDs (\texttt{sp\_*}, \texttt{m\_*}, \texttt{rq\_*}) leaked into prose --- checked directly against the discussion pipeline log.
\item \textbf{Notation with worked example} --- the formal setup section defines $(D_N, \alpha, E, \hat\alpha, g, B, e_0(N), e_B(N), \varepsilon, \beta)$ and the Problem Definition/Methodology sections consistently reuse this notation; no separate worked numeric example is included since no numbers exist yet (would require fabricating results).
\item \textbf{RQ framing matches outcomes} --- all 7 RQs are graded/quantitative (rates, ratios, confidence-interval bounds), each answered in its own subsection with an explicit falsifiable expected outcome \emph{and} a distinct falsification condition, not a binary yes/no.
\item \textbf{No typesetting artifacts} --- only explicit placeholder markers in Results, per the drafting convention this is a pre-registered spec, not a report.
\item \textbf{Artifact plan} --- the estimator/adversary harness and the self-generated semi-synthetic corpora are the natural release artifacts; not yet built, so not claimed as available.
\end{itemize}

\section*{Hand-off}
This is a draft, not a finished paper --- experiments have not run. Natural next steps: prose polish and venue-specific LaTeX formatting once a target venue is chosen, or revising an earlier stage if piloting surfaces a problem (e.g., if the shuffled-label control in RQ4 fails).

\bibliographystyle{plainnat}
\bibliography{refs}

\end{document}
